% 题证摘录；来源与整理边界见相同编号分题文件。

\paragraph{2. Moser's Theorem.}
One can ask whether, for a given manifold $M$, two symplectic structures $\omega_0,\omega_1$ are equivalent, i.e. whether there is a symplectomorphism $M\to M$ which pulls back one to the other. In general, $[\omega_0]=[\omega_1]$ does not imply that the two structures are symplectomorphic. To study this question further, we give other notions of equivalence.
\begin{definition}[2]
Two forms $\omega_0,\omega_1$ are deformation equivalent if $\exists(\omega_t)_{t\in[0,1]}$ a continuous family of symplectic forms, and isotopic if there is such a family with $[\omega_t]$ constant in $H^2(M,\R)$.
\end{definition}
Let $M$ be a compact manifold with $\omega_0,\omega_1$ isotopic symplectic forms (i.e. $\exists\omega_t$ as above with each $\omega_t$ nondegenerate).
\begin{theorem}[1: Moser]
$\exists$ an isotopy $\rho_t:M\to M$ s.t. $\rho_t^*\omega_t=\omega_0$.
\end{theorem}
That is, $(M,\omega_0)$ and $(M,\omega_1)$ are symplectomorphic.
\begin{proof}
(This technique is known as Moser's trick.) By assumption, $[\omega_t]$ is independent of $t$, i.e. $[\frac{d\omega_t}{dt}]=0$. Thus, $\exists\alpha_t$ a 1-form s.t. $\frac{d\omega_t}{dt}=-d\alpha_t$: moreover, we can choose this $\alpha_t$ smoothly w.r.t. to $t$ (via the Poincaré lemma). Since $\omega_t$ is nondegenerate, $\exists X_t$ s.t. $i_{X_t}\omega_t=\alpha_t$. Moreover, since $M$ is compact, we have a well-defined flow $\rho_t$ of $X_t$. Now,
\begin{align*}
\frac{d}{dt}(\rho_t^*\omega_t)
&=\rho_t^*(L_{X_t}\omega_t)+\rho_t^*\left(\frac{d\omega_t}{dt}\right)\\
&=\rho_t^*\left(di_{X_t}\omega_t+\frac{d\omega_t}{dt}\right)=0.
\end{align*}
Since $\rho_0$ is the identity, we have our desired isotopy.
\end{proof}

\begin{theorem}[2: Darboux]
For $(M,\omega)$ symplectic, $p\in M$, $\exists U\ni p$ with a coordinate system $(x_1,y_1,\ldots,x_n,y_n)$ s.t. $\omega|_U=\sum dx_i\wedge dy_i$.
\end{theorem}
\begin{proof}
$(T_pM,\omega_p)$ has a standard basis $(e_1,\ldots,e_n,f_1,\ldots,f_n)$, so there exist local coordinates $(x_1,y_1,\ldots,x_n,y_n)$ s.t. $\omega_p=\sum dx_i\wedge dy_i$. On a neighborhood $U$ of $p$, we obtain two symplectic forms: $\omega$ and the standard form. The family $\omega_t=(1-t)\omega_0+t\omega_1$ is one of closed forms: since nondegeneracy is an open condition, we can shrink our neighborhood to assure that $\omega_t$ is nondegenerate for each $t$ on some $U'\ni p$. Thus, $\exists\alpha\in\Omega^1(U)$ s.t. $\omega_1-\omega_0=-d\alpha$. Subtracting a constant, we can assume $\alpha_p=0$. Let $v_t$ be the vector field on $U$ s.t. $i_{v_t}\omega_t=\alpha$. Then $\exists U''\ni p$ s.t. its flow $\rho_t$ is defined $\forall t$. By the Moser's trick, we fnd that $\rho_1^*\omega_1=\omega_0$, implying that the symplectic form is indeed standard after composing our chosen coordinates with $\rho_1$.
\end{proof}
